\documentclass[11pt]{article}
\usepackage[margin=1in]{geometry}
\usepackage{amsmath,amssymb}
\usepackage{booktabs}
\usepackage{array}
\usepackage{xcolor}
\usepackage[colorlinks=true,linkcolor=blue,citecolor=blue,urlcolor=blue]{hyperref}
\usepackage{enumitem}

\newcommand{\B}[1]{\textbf{#1}}
\newcommand{\cfauc}{\mbox{cf-AUC}}

\title{\bf Isolating why future-reachability values fail on cube manipulation\\[3pt]
\large An FB\,$\leftrightarrow$\,GCIQL bridge --- experiment log}
\author{Factored-FB}
\date{2026-06-29 \quad(\emph{living document}; two tiers --- Part~1 development
(500k, 3 seeds; the goal-matching ladder) and Part~2 final (1M, 10 seeds). Spine
revised: FB-IQL critic is $\sim$59 at sane lr/ortho, not value-blind at $24$.)}

\begin{document}
\maketitle

\section{Starting point: the updated performance table}

All methods use the FlowBC actor (the only actor we study); each is its tuned
configuration after a hyperparameter study. \B{The table below reports each method's
best eval over training} (peak IQM over steps $\ge$500k --- most methods peak before
1M and several degrade after); \B{the later ablation sections report success at 1M}
($n{=}8$--$10$ seeds), where the FB-side interventions and G4 live. ($\pm$ is a $95\%$ CI, $1.96\times$ the
bootstrap standard error over $n{=}10$ seeds; success is IQM over seeds, in \%, from
wandb \texttt{amsks/factored-fb}.)

\begin{center}
\begin{tabular}{lc}
\toprule
Method (state, FlowBC, best $\ge$500k) & IQM success ($\pm$\,95\% CI) \\
\midrule
GCIQL & $78 \pm 4$ \\
FB    & $52 \pm 6$ \\
PSM   & $42 \pm 8$ \\
CRL   & $31 \pm 3$ \\
RLDP  & $27 \pm 6$ \\
\bottomrule
\end{tabular}
\end{center}
\noindent (Peak step per method.)

\noindent Every future-reachability method (FB / PSM / CRL / RLDP) sits far below
GCIQL. The rest of this log is the set of experiments we ran to find \emph{which}
ingredient of the FB family causes that gap.

\section{The two objectives, and the modifications they motivate}

\subsection{FB objective}
FB represents the discounted successor measure of a \emph{policy family} $\{\pi_z\}$.
With data distribution $\rho$ and a task vector $z$,
\[
M^{\pi_z}(s,a,\mathrm{d}s') := \sum_{t\ge 0}\gamma^{t}\,
\Pr\!\big(s_{t+1}\in\mathrm{d}s'\mid s_0{=}s,a_0{=}a,\pi_z\big)
\;\approx\; F(s,a,z)^\top B(s')\,\rho(\mathrm{d}s').
\]
$F,B$ are fit by the residual of the measure Bellman equation
$M = P + \gamma P^{\pi_z} M$ --- a TD over \emph{random} goals $s'\sim\rho$ plus a
diagonal term, under the orthonormality constraint $\mathbb{E}_\rho[BB^\top]=I$
(coefficient $\lambda$ $=$ our \texttt{ortho\_coef}); schematically,
\begin{align}
\mathcal{L}_{\mathrm{FB}}
&= \mathbb{E}_{\substack{(s,a,s^{+})\sim\rho\\ s'\sim\rho,\ z}}
\Big[\big(F(s,a,z)^\top B(s')
   - \gamma\,\mathrm{sg}\big[\bar F(s^{+},\pi_z(s^{+}),z)^\top \bar B(s')\big]\big)^2\Big]
\nonumber\\[-2pt]
&\quad\;-\;2\,\mathbb{E}_{(s,a,s^{+})\sim\rho,\ z}\big[F(s,a,z)^\top B(s^{+})\big]
\;+\;\lambda\,\big\|\mathbb{E}_\rho[BB^\top]-I\big\|_F^2 .
\end{align}
The reward never enters training. A task is specified \emph{post hoc} by a linear
readout, and the policy is greedy on a bilinear $Q$:
\[
z_r=\mathbb{E}_{s'\sim\rho}\big[r(s')\,B(s')\big],\qquad
Q_z(s,a)=F(s,a,z)^\top z,\qquad
\pi_z(s)=\arg\max_a Q_z(s,a).
\]

\subsection{GCIQL objective}
GCIQL learns a \emph{goal-conditioned optimal} value by Implicit Q-Learning on the
sparse goal-reaching reward $r(s,g)=\mathbb{1}[s\!\approx\!g]-1\in\{-1,0\}$
(episode ends at the goal; $d=\mathbb{1}[s\!\approx\!g]$). An expectile value
approximates the in-support $\max_a$, and $Q$ bootstraps on it:
\begin{align}
\mathcal{L}_V &= \mathbb{E}_{(s,a,g)\sim\rho}\big[L_2^{\tau}\!\big(\bar Q(s,a,g)-V(s,g)\big)\big],
\qquad L_2^{\tau}(u)=\big|\tau-\mathbb{1}[u<0]\big|\,u^2,\ \ \tau>\tfrac12,\\
\mathcal{L}_Q &= \mathbb{E}_{(s,a,s^{+},g)\sim\rho}\Big[\big(\,
\underbrace{r(s,g)+\gamma(1-d)\,V(s^{+},g)}_{=\;(\mathrm{succ}-1)+\gamma(1-\mathrm{succ})\,V(s^{+},g)}
\;-\;Q(s,a,g)\big)^2\Big].
\end{align}
The goal $g$ is a \emph{network input}; the reward enters the backup at \emph{train}
time; the expectile $L_2^{\tau}$ is an implicit policy-improvement (greedy) operator.
The underbrace is exactly the bug/fix from the FB-IQL experiment below: masking the
bootstrap by goal-reached, not by the env \texttt{terminated} flag.

\subsection{Key differences}
\begin{enumerate}[nosep]
\item \B{Backup signal --- density vs reward.} FB regresses the \emph{occupancy
  measure} (every $s'\sim\rho$ is a goal at once); GCIQL regresses the \emph{value} of
  one goal's reward. ($\bar F^\top\bar B$ over random $s'$ vs $r+\gamma V$ for $g$.)
\item \B{Improvement operator --- none vs expectile-max.} FB's measure backup is
  \emph{policy evaluation} of $\{\pi_z\}$ --- no $\max$ appears in the residual.
  GCIQL's $L_2^{\tau}$ drives $V$ to an upper expectile of $Q$, an \emph{implicit
  $\max_a$}, so the backup is greedy/optimal.
\item \B{Task entry --- post-hoc linear $z$ vs train-time input $g$.} FB infers
  $z=\mathbb{E}[rB]$ at test; GCIQL conditions on $g$ throughout training.
\end{enumerate}

\subsection{The modifications (each isolates one difference)}

\paragraph{M1 --- objective: measure $\to$ goal-conditioned value-TD.}
Replace $\mathcal{L}_{\mathrm{FB}}$ by $\mathcal{L}_V+\mathcal{L}_Q$ on FB's \emph{own}
bilinear $Q=F\!\cdot\!z$ (the \emph{FB-IQL critic}). This injects both the reward
backup and the expectile max (differences 1\,\&\,2) while holding the form and
inference fixed:
\[
\mathcal{L}_{\mathrm{FB}}\ \longmapsto\
\mathbb{E}\big[L_2^{\tau}(\bar Q-V)\big]
+\mathbb{E}\big[(r+\gamma(1-d)V' - Q)^2\big],\qquad Q=F(s,a,z)^\top z .
\]
\emph{Result: 24 --- a correct value on the measure-trained bilinear stays
value-blind.}

\paragraph{M2 --- form: bilinear, IQL-trained, exact goal.}
Keep the inner-product parameterization but train it by M1's objective with the exact
goal as input:
\[
Q(s,a,g)=\frac{\phi(s,a)^\top\psi(g)}{\sqrt d},\qquad
V(s,g)=\frac{\phi_V(s)^\top\psi_V(g)}{\sqrt d},
\]
plugged into $\mathcal{L}_V,\mathcal{L}_Q$. This is the \emph{bilinear bridge}: FB's
form and rank, GCIQL's objective and conditioning. \emph{Result: 94.}

\paragraph{M3 --- conditioning: $z$-condition the actor under FB's inference.}
On the IQL-trained value of M2, set $B:=\psi$ and condition the policy on a task
vector rather than the raw goal,
\[
z=\psi(g)\ \ (\text{exact})\qquad\text{or}\qquad
z=\mathbb{E}_{s'\sim\rho}\big[r(s')\,\psi(s')\big]\ \ (\text{FB's L0 rule}),
\qquad \pi(\cdot\mid s,z),\ \ Q=\phi(s,a)^\top z .
\]
For a sparse single-goal reward $r(s')=\mathbb{1}[s'\!\approx\!g]$ the readout collapses
to $z=\mathbb{E}[r\,\psi(s')]\propto\psi(g)$, exactly recovering FB's linear inference
--- but now $\psi$ is IQL-trained, not measure-trained. This is \emph{G4}: the
success-level test of whether the \emph{objective alone}, under FB's full zero-shot
pipeline, suffices (\S\,G4 below: yes with exact $z$, no with inferred $z$).

\medskip
\noindent In short: the FB$\to$GCIQL gap is three structural differences (signal,
improvement operator, task entry); M1 changes the first two together (the decisive
one), M2 the form they ride on, M3 the third. The experiments below are these three
modifications and the controls around them.

\section{Two reporting tiers}
\noindent The results split into two tiers; read them apart. \textbf{Part~1 --- development
(500k steps, 3 seeds $\{0,1,2\}$):} the fast tier where we \emph{figure things out}; the
headline is the goal-matching ladder, and conclusions are provisional ($n{=}3$, wide CIs).
\textbf{Part~2 --- final (1M steps, 10 seeds):} the slow, confirmed tier (the full ablation
grid + the elimination experiments at scale). \emph{Where the two disagree, Part~1 is the
current best evidence and Part~2 is being re-run at the corrected configuration} --- most
importantly the FB-IQL critic (Part~2 grid: $24$; Part~1 ladder at sane lr/ortho: $\sim$59;
see the Revision note in ``Importing the value objective\ldots'' in Part~2).

\section{Part~1 --- Development (500k, 3 seeds): the goal-matching ladder}
The fast-iteration tier --- the full FB$\to$GCIQL morph at \textbf{500k} steps,
3 seeds $\{0,1,2\}$, the \emph{flow actor held fixed across every row}, with the
GCIQL \emph{dataset goal-matching} reward (success $=$ goal is the current state).
Critic and actor use \emph{separate} goal sampling: critic mix
$[p_{\text{traj}},p_{\text{cur}},p_{\text{rand}}]=[0.5,0.2,0.3]$ (the $0.2$ cur-goals
are the reached anchors that pin $V$ to $0$ at the goal), actor trajectory-only
$[1,0,0]$ --- exactly GCIQL's value/actor split. $\pm$ is the 95\% CI (Student-$t$, $n{=}3$).
\begin{center}\small
\begin{tabular}{lc}
\toprule
Method / step (all FlowBC actor, @500k) & success ($\pm$\,95\% CI) \\
\midrule
\multicolumn{2}{l}{\emph{Future-reachability baselines}}\\
CRL (flowbc, contrastive) & $30 \pm 5$ \\
plain FB (measure) & $34 \pm 15$ \\
\midrule
\multicolumn{2}{l}{\emph{FB-IQL (objective swapped: measure $\to$ goal-matching value)}}\\
FB-IQL, goal-match $+$ GCIQL mix & $\mathbf{59 \pm 15}$ \\
\quad $+$ de-condition $F\!\to\!\phi$ & $58 \pm 26$ \\
\quad $+$ $\sqrt d$ normalization & $65 \pm 11$ \\
\quad $+$ both & $57 \pm 4$ \\
\midrule
\multicolumn{2}{l}{\emph{GCIQL targets}}\\
stock GCIQL $+$ flow & $71 \pm 6$ \\
bilinear-GCIQL ($d{=}50$) & $\mathbf{90 \pm 6}$ \\
\bottomrule
\end{tabular}
\end{center}
\noindent\textbf{Reading it.} The \emph{objective} is the lever: swapping FB's measure
backup for the IQL goal-matching value nearly \emph{doubles} success (CRL/plain FB
$\sim$30--34 $\to$ FB-IQL $59$), with the flow actor and the bilinear/learned-$B$ form
held fixed. The \emph{structural} knobs --- de-conditioning the forward map ($F\!\to\!\phi$)
and the $\sqrt d$ scaling --- are \emph{flat} (57--65, all overlapping at $n{=}3$). A real
gap to GCIQL remains ($59\to71\to90$) that these knobs do not close (candidates: a bilinear
$V$ head, the clipped-double-$Q$ minimum, a learned $\psi$). \textbf{This revises the
$\S$\,``importing the value objective doesn't help'' result in Part~2:} that section's FB-IQL
critic $=24$ used $\mathrm{lr}{=}1e\text{-}4$ and ortho$=1000$ and collapsed on 2/3 seeds;
at $\mathrm{lr}{=}3e\text{-}4$, ortho$=0$, and the dataset goal-match reward, FB-IQL is
\emph{stable at} $\sim$59 --- so the original $24$ was substantially an lr/ortho
training-stability artifact, not evidence the bilinear value is intrinsically value-blind.
\emph{Caveats:} $n{=}3$ @500k; plain FB is still climbing at 500k (its @1M $=48$); lr and
ortho are not yet isolated from each other.

\section{Part~2 --- Final (1M, 10 seeds): the morph bridge at scale}
We use a two-directional ``morph bridge'': start from FB and import GCIQL's value
objective; and start from GCIQL and impose FB's restrictions one at a time. Each
subsection is one experiment: what we changed, and what success it produced. \emph{These
are the 1M/10-seed results; the Part~1 ladder supersedes any that conflict (FB-IQL critic).}

\subsection{The actor is not the cause}
\textbf{Did:} held the FlowBC actor \emph{constant} and varied only the critic it
steers on. The same flow actor sits on FB's measure-trained representation, on
GCIQL's MLP critic, and on the IQL-trained bilinear critic (\S3.2). \\
\textbf{Found:} with one fixed actor, success spans \B{$48\pm5$} (FB rep) $\to$
\B{$73\pm4$} (GCIQL MLP critic) $\to$ \B{$94\pm1$} (IQL bilinear). \\
\textbf{Implies:} the FB-family deficit lives in the \emph{value path}
(representation $\times$ readout $\times$ inference), not the actor --- the actor is
identical across the whole range. This is the constant baseline for all experiments
below.

\subsection{The bilinear form is not the cause (the bridge)}
\textbf{Did:} replaced GCIQL's MLP critic with an FB-shaped \emph{bilinear} value
$Q=\phi(s,a)^\top\psi(g)/\sqrt{d}$ --- the same factorization FB/PSM use --- still
IQL-trained, still handed the exact goal, with the flow actor. (Group
\texttt{GCIQL-State-Flow-Bilinear}, $d{=}512$.) \\
\textbf{Found:} success \B{$94\pm1$} ($n{=}8$) --- \emph{above} the MLP critic
($73\pm4$), and an FB-shaped value essentially solves the task. \\
\textbf{Implies:} the bilinear factorization $Q=F\cdot z$ is \B{not} the wall. An
inner-product value can carry full controllability when trained the right way.

\subsection{Rank is not the cause (down to a far-below-FB floor)}
\textbf{Did:} re-ran the bilinear bridge across ranks $d\in\{5,50,128,512\}$. \\
\textbf{Found:} $d{=}50$: $\mathbf{93\pm1}$; $d{=}128$: $\mathbf{94\pm1}$;
$d{=}512$: $\mathbf{94\pm1}$ --- \emph{flat} across this $10\times$ range; but
$d{=}5$: $\mathbf{14\pm3}$ --- a \emph{collapse} (already $\sim5\%$ at 100k
vs.\ $\sim70\%$ for $d{\ge}50$). \\
\textbf{Implies:} there is a capacity cliff between $d{=}5$ and $d{=}50$, but it sits
\B{far below FB's operating range} ($d{=}50$--$128$). So rank is \B{not} FB's crutch:
at FB-parity rank the IQL-trained bilinear already scores $\sim93$, and only an
extreme $d{=}5$ (where the inner-product can no longer represent the goal-reaching
structure) breaks it.

\subsection{Full ablation grid}
The complete set of single-knob modifications, starting from stock GCIQL and morphing
one ingredient at a time toward FB's bilinear measure-matching form. All numbers are
success at \textbf{1M} (state, FlowBC actor unless noted). We report the grid at
\emph{two tiers}: a \textbf{development set} (Grid~A) of the \emph{same three
preliminary seeds} (0,1,2) used to scope each ablation, and the \textbf{final}
10-seed campaign (Grid~B). The development tier is where new ablations enter --- a row
is run on seeds 0--2 first and added to Grid~A, then graduated to Grid~B once the
remaining seeds (3--9) complete. Grid~A's seeds are a \emph{strict subset} of Grid~B's,
so the two tiers are directly comparable; Grid~B simply tightens every interval
(e.g.\ single-$Q$ $90.5\pm4.0$ at $n{=}3 \to 90.2\pm1.5$ at $n{=}10$). $\pm$ is the
\textbf{95\% CI} (Student-$t$, $n{-}1$ df; $t{=}4.30$ at $n{=}3$, so the 3-seed
intervals are wide by design).

\paragraph{Grid~A --- development set (3 seeds @1M).} Every ablation on seeds
$\{0,1,2\}$, the apples-to-apples tier for incoming ablations.
\begin{center}
\small
\begin{tabular}{lcc}
\toprule
\textbf{Modification} & \textbf{Success @1M (95\% CI)} & \textbf{$n$} \\
\midrule
\multicolumn{3}{l}{\emph{Baselines (MLP critic)}}\\
\textit{none} --- stock GCIQL (MLP critic, Gaussian DDPG+BC actor) & ---$^{\ddagger}$ & --- \\
actor $\to$ \textbf{flow (FQL)} actor (``GCIQL--FlowBC'')          & $73.6 \pm 8.1$ & 3 \\
\midrule
\multicolumn{3}{l}{\emph{Critic form --- the key lever}}\\
critic $\to$ \textbf{bilinear} $Q{=}\phi\!\cdot\!\psi/\sqrt d$, $d{=}50$ ($+$flow) & $\mathbf{91.6 \pm 5.1}$ & 3 \\
\midrule
\multicolumn{3}{l}{\emph{GCIQL design knobs (on bilinear $d{=}50$ $+$ flow)}}\\
\textbf{Gaussian} DDPG+BC actor (no flow)            & $56.1 \pm 15.1$ & 3 \\
expectile $\tau{=}0.5$ (no pessimism)                & $88.7 \pm 1.7$  & 3 \\
expectile $\tau{=}0.95$ (strong pessimism)           & $90.3 \pm 3.6$  & 3 \\
\textbf{single-}$Q$ (no double-$Q$ minimum)          & $90.5 \pm 4.0$ & 3 \\
actor on $z{=}\psi(g)$, not raw goal (G4 exact-$z$)   & $92.4 \pm 3.8$ & 3 \\
\midrule
\multicolumn{3}{l}{\emph{Rank sweep (bilinear $+$ flow)}}\\
rank $d{=}5$    & $13.3 \pm 7.1$  & 3 \\
rank $d{=}50$   & $91.6 \pm 5.1$ & 3 \\
rank $d{=}128$  & $93.1 \pm 1.1$ & 3 \\
rank $d{=}512$  & $94.3 \pm 3.0$ & 3 \\
\midrule
\multicolumn{3}{l}{\emph{FB objective --- the bridge back}}\\
plain FB --- bilinear $+$ flow, measure-matching, zero-shot $z$ & $48.0 \pm 5.0$ & 3 \\
FB, measure-matching, trajectory-$z$ (zHER)         & $29.3 \pm 18.5$ & 3 \\
FB--IQL, goal-cond.\ IQL critic \emph{replaces} measure loss & $\mathbf{18.0 \pm 26.5}$ & 3 \\
FB--IQL, goal-TD $q$-loss added (measure kept, $q\_loss\_coef{=}1$) & $30.7 \pm 16.6$ & 3 \\
\midrule
\multicolumn{3}{l}{\emph{CRL (contrastive future-reachability baseline)}}\\
CRL, flow actor               & $27.6 \pm 4.6$ & 3 \\
CRL, Gaussian DDPG+BC actor   & $15.1 \pm 3.1$ & 3 \\
\bottomrule
\end{tabular}
\end{center}
\noindent\footnotesize $^{\ddagger}$ Stock GCIQL is a 12-seed reuse from an earlier
campaign; its per-seed values for the $\{0,1,2\}$ subset were not relocated, so only the
Grid~B aggregate ($67.4$) is available. G4 inferred-$z$ ($6.7$) is a dedicated 50-episode
zero-shot eval and lives only in Grid~B. \normalsize

\paragraph{Grid~B --- final (10 seeds @1M).} The full campaign (campaign spec
\texttt{2026-06-22-bridge-table-10seed-1M}), $n{=}8$--$10$ seeds per row; each block is
expanded in the surrounding sections.
\begin{center}
\small
\begin{tabular}{lcc}
\toprule
\textbf{Modification} & \textbf{Success @1M (95\% CI)} & \textbf{$n$} \\
\midrule
\multicolumn{3}{l}{\emph{Baselines (MLP critic)}}\\
\textit{none} --- stock GCIQL (MLP critic, Gaussian DDPG+BC actor) & $67.4 \pm 4.6$ & 10 \\
actor $\to$ \textbf{flow (FQL)} actor (``GCIQL--FlowBC'')          & $72.8 \pm 4.3$ & 9 \\
\midrule
\multicolumn{3}{l}{\emph{Critic form --- the key lever}}\\
critic $\to$ \textbf{bilinear} $Q{=}\phi\!\cdot\!\psi/\sqrt d$, $d{=}50$ ($+$flow) & $\mathbf{92.6 \pm 1.3}$ & 10 \\
\midrule
\multicolumn{3}{l}{\emph{GCIQL design knobs (on bilinear $d{=}50$ $+$ flow)}}\\
\textbf{Gaussian} DDPG+BC actor (no flow)            & $56.6 \pm 5.3$ & 10 \\
expectile $\tau{=}0.5$ (no pessimism)                & $88.9 \pm 2.3$  & 10 \\
expectile $\tau{=}0.95$ (strong pessimism)           & $90.6 \pm 1.2$  & 10 \\
\textbf{single-}$Q$ (no double-$Q$ minimum)          & $90.2 \pm 1.5$ & 10 \\
actor on $z{=}\psi(g)$, not raw goal (G4 exact-$z$)$^{\dagger}$ & $93.0 \pm 1.9$ & 9 \\
\quad same, zero-shot $z{=}\mathbb{E}[r\psi(s')]$ (G4 inferred-$z$) & $6.7 \pm 3.1$ & 9 \\
\midrule
\multicolumn{3}{l}{\emph{Rank sweep (bilinear $+$ flow)}}\\
rank $d{=}5$    & $14.3 \pm 3.0$  & 10 \\
rank $d{=}50$   & $92.6 \pm 1.3$ & 10 \\
rank $d{=}128$  & $93.9 \pm 1.3$ & 10 \\
rank $d{=}512$  & $94.5 \pm 1.3$ & 8 \\
\midrule
\multicolumn{3}{l}{\emph{FB objective --- the bridge back}}\\
plain FB --- bilinear $+$ flow, measure-matching, zero-shot $z$ & $48 \pm 5$ & 10 \\
FB, measure-matching, trajectory-$z$ (zHER)         & $32.0 \pm 5.3$ & 10 \\
FB--IQL, goal-cond.\ IQL critic \emph{replaces} measure loss & $\mathbf{23.6 \pm 6.8}$ & 10 \\
FB--IQL, goal-TD $q$-loss added (measure kept, $q\_loss\_coef{=}1$) & $38.8 \pm 6.2$ & 10 \\
\midrule
\multicolumn{3}{l}{\emph{CRL (contrastive future-reachability baseline)}}\\
CRL, flow actor               & $25.8 \pm 2.4$ & 9 \\
CRL, Gaussian DDPG+BC actor   & $16.8 \pm 2.7$ & 10 \\
\bottomrule
\end{tabular}
\end{center}
\noindent\footnotesize $^{\dagger}$ G4 row is the dedicated 50-episode zero-shot
evaluation (\S\,G4), $n{=}9$ seeds at 1M; exact-$z$ $93.0\pm1.9$, inferred-$z$
$6.7\pm3.1$. At 1M the dedicated exact-$z$ eval agrees with the in-training eval
($93.4$), unlike at 500k where the vendored loop's raw-goal feed made them diverge.
$\pm$ is the 95\% CI (Student-$t$, $n{-}1$ df); the per-seed values are in the
verification log. \normalsize

\smallskip
\noindent\textbf{Reading the grid.} The bridge's $\sim$93 is \emph{robust} to GCIQL's
pessimism knobs --- removing the expectile max ($\tau{=}0.5$, $89$) or the double-$Q$
minimum ($90$) barely moves it. The one design choice that matters is the
\emph{actor}: a Gaussian DDPG+BC actor on the \emph{same} bilinear value drops to $57$,
and the same flow$\to$Gaussian swap costs CRL $\sim9$ points ($26\!\to\!17$) --- so the
flow actor is needed to \emph{realize} a good bilinear value, but it is held
\emph{fixed} across every FB-vs-GCIQL comparison (\S3.4) and does not explain the gap.
CRL, a future-reachability method like FB/PSM/RLDP, sits with them far below the bridge
under \emph{either} actor.

\subsection{Zero-shot inference is not the cause (value level)}
\textbf{Did:} on the bilinear bridge value, ranked goal-reaching states by
$V_z(s)=\phi(s)^\top z$ using (i) the exact goal embedding $z{=}\psi(g)$ vs.\ (ii)
FB's own zero-shot rule $z{=}\mathbb{E}[r\,\psi(s')]$. We measured rank-AUC of the
value over reached vs.\ not-reached states (\cfauc{}; scale-invariant in $z$, so no
sphere-projection ambiguity). 4 seeds. \\
\textbf{Found:} \cfauc{} \B{$0.922{\pm}.020$ (exact) vs.\ $0.915{\pm}.019$
(inferred)} --- identical. The same inference rule leaves FB's own value near chance
($\approx0.5$). \\
\textbf{Implies:} FB-style zero-shot inference is \B{not} the differentiator. The
IQL-trained value ranks goal-reaching states equally well under FB's own inference;
the difference must be upstream, in the representation the objective builds.

\subsection{Importing the value objective into FB's framework does not help}
\textbf{\textcolor{red}{Revision (see the 500k development ladder above).}} The headline of
this section --- ``FB-IQL critic $=24$, value-blind'' --- is \textbf{superseded}. Those runs
used $\mathrm{lr}{=}1e\text{-}4$, ortho$=1000$, and the privileged cube-xyz reward, and collapsed
on 2/3 seeds. With $\mathrm{lr}{=}3e\text{-}4$, ortho$=0$, and the GCIQL dataset goal-matching
reward, the \emph{same} IQL-critic objective is \textbf{stable at $\sim$59} (@500k, $n{=}3$) ---
\emph{above} plain FB, doubling the future-reachability baselines. So the $24$ was largely an
lr/ortho training-stability artifact, \emph{not} evidence the bilinear value is intrinsically
value-blind. The corrected reading: the \emph{objective} (measure $\to$ goal-matching value-TD)
\emph{is} the lever; it just stops short of GCIQL ($59 \to 90$), and the structural knobs we
tried (de-condition $F\!\to\!\phi$, $\sqrt d$) do not close that residual. The $@1M$ numbers in
the table below are retained as the \emph{old} measure-framework interventions (Part~2 will
re-confirm them at the corrected incidentals).

\medskip
\textbf{Did:} four ways to give FB the value objective while keeping its
bilinear/measure framework:
\begin{center}
\begin{tabular}{lc}
\toprule
Intervention & success ($\pm$\,95\% CI) \\
\midrule
plain FB (anchor, $n{=}10$) & $48 \pm 5$ \\
$q\_loss\_coef{=}1$ (up-weight goal-TD on $Q{=}F\!\cdot\!z$) & $39 \pm 6$ \\
zHER (trajectory-future task sampling) & $32 \pm 5$ \\
IWR (interaction-weighted measure, $n{=}3$) & $23 \pm 23$ \\
FB-IQL critic (measure loss \B{replaced} by goal-cond.\ IQL) & $\mathbf{24 \pm 7}$ \\
decouple actor-$z$ (reachable-$z$ actor, perm critic goals; $n{=}3$) & $39 \pm 34$ \\
decouple actor-$z$ (reachable-$z$ actor, traj critic goals; $n{=}3$) & $18 \pm 13$ \\
\midrule
\multicolumn{2}{l}{\emph{+ drop the learned backward map ($z{=}g$, $B{=}I$, on-traj $z$; $n{=}3$, @1M):}}\\
\quad FB-IQL critic, $z{=}g$ & $\mathbf{44 \pm 13}$ \\
\quad zHER, $z{=}g$ (measure loss; confounded) & $19 \pm 22$ \\
\midrule
\multicolumn{2}{l}{\emph{single-variable controls for the $z{=}g$ rows ($n{=}3$, @1M):}}\\
\quad FB-IQL critic, learned $B$, ratio${=}1$ & $21 \pm 40$ \\
\quad zHER, $z{=}g$, $z$ re-normalized & $20 \pm 5$ \\
\bottomrule
\end{tabular}
\end{center}
\noindent (interventions $n{=}10$ except IWR, the decouple rows, the $z{=}g$ rows, and
the controls, all $n{=}3$.)
\textbf{Found:} every intervention \emph{that keeps the learned backward map} lands
\emph{below} plain FB (the $z{=}g$ rows are the exception, below). The FB-IQL critic is
the sharpest: a first version diverged (V/Q ramped to the $1/(1{-}\gamma){=}100$
fixed point, actor collapsed to $3$) because the Bellman bootstrap was masked by the
env \texttt{terminated} flag (never set in cube offline data) instead of by
goal-reached. With the GCIQL-faithful target
$y=(\text{succ}{-}1)+\gamma(1{-}\text{succ})V'$ the value becomes bounded and stable
--- and success is still only \B{24}. \\
\textbf{Implies:} a \emph{correctly-trained} goal-conditioned IQL value on FB's
\emph{measure-trained} bilinear is value-blind. The value-target form is not the
lever; the representation the measure objective builds is the problem.\\
\textbf{Dropping the learned backward map ($z{=}g$, $n{=}3$, @1M).} Pushing further, we
\emph{drop the learned backward map entirely} and feed the goal directly as the task
vector ($B{=}I$, $z{=}g$), with all task sampling on-trajectory (\texttt{train\_goal\_ratio}${=}1$).
With the goal-conditioned IQL critic this recovers to \B{$44\pm13$} --- statistically
indistinguishable from plain FB ($48\pm5$) and well above the learned-$B$ FB-IQL critic
($24$). So the FB-IQL critic's collapse is substantially the \emph{bilinear backward-map
bottleneck}, not the value target: remove $B$ and the value objective is no longer
value-blind. Crucially this only \emph{recovers the FB ceiling} ($\sim$48); it does not
approach GCIQL ($\sim$90), so the broader gap stands. The measure-loss variant with
$z{=}g$ instead \emph{degrades} ($19\pm22$) and is confounded --- with $B{=}I$ the
measure's implied reward $z^\top B(s){=}g^\top s$ is linear, not goal-reaching, and the
successor measure diverges.

\textbf{The two single-variable controls (now run) isolate the backward map.} The $z{=}g$
intervention changes two things at once relative to the FB-IQL critic row (drops $B$ \emph{and}
moves task sampling to on-trajectory at ratio${=}1$). The controls split them. (i) Learned-$B$
FB-IQL critic at \texttt{train\_goal\_ratio}${=}1$ scores \B{$21\pm40$} --- essentially the same
as the ratio$\,{=}0.5$ FB-IQL critic ($24$), so the sampling ratio is \emph{not} the lever; the
$24\!\to\!44$ jump is attributable to dropping $B$. (Indeed, under \texttt{critic\_mode=iql} both
the decouple-actor-$z$ knob and \texttt{train\_goal\_ratio} are inert: the learned-$B$ ratio$\,{=}1$
run is byte-identical to the decouple-IQL run across all eval checkpoints, confirming they collapse
to one configuration.) (ii) Re-normalizing $z$ in the $z{=}g$ measure variant gives \B{$20\pm5$},
no better than the un-normalized $19\pm22$ --- bounding $z$ does not rescue the measure blow-up, so
that row's weakness is the linear implied reward, not an unbounded-$z$ artifact. Together: the
learned backward map $B$ is the bottleneck the $z{=}g$ rows expose, cleanly separated from the
sampling and normalization changes.

\textbf{Decoupling the actor's $z$ from the critic's also does not help.} A natural in-framework fix
is to let the critic keep its broad task distribution (random directions $+$ off-diagonal negatives)
while the \emph{actor} trains only on reachable tasks $z{=}B(\text{on-traj future})$ --- mirroring
GCIQL's separation of what the value learns about from what the actor practices. In measure mode this
lands at \B{$39\pm34$} (perm critic goals) and \B{$18\pm13$} (trajectory critic goals), both below
plain FB ($48$); under the IQL critic the knob is inert (above). So decoupling the actor's task
distribution does not recover controllability either --- consistent with the rest of this table:
nothing that keeps the learned backward map and the measure framework lifts FB above its own ceiling.

\subsection{``Improving the representation'' backfires}
\textbf{Did:} measured how well each method's backward map / basis encodes the
interaction variable (nonlinear lift-$R^2$, separability) and compared to success;
plus the PSM/RLDP HP sweeps and the $d{=}128$ FB run. \\
\textbf{Found:} success \B{anti-correlates} with interaction-encoding:
\begin{center}
\begin{tabular}{lccc}
\toprule
Config & basis lift-$R^2$ & separability & success ($\pm$\,95\% CI) \\
\midrule
plain FB ($d{=}50$) & 0.02 & 0.74 & $\mathbf{48 \pm 5}$ \\
FB ($d{=}128$)      & 0.30 & 0.905 & $28 \pm 4$ \\
PSM (ortho1000)     & 0.81 & --   & $35 \pm 4$ \\
IWR ($d{=}50$)      & 0.79 & 0.89 & $23 \pm 23$ \\
\bottomrule
\end{tabular}
\end{center}
IWR does exactly what its derivation predicts --- it re-orders the spectrum so the
basis concentrates on the interaction variable (lift $0.02\!\to\!0.79$) --- and
success \emph{drops}. The training-task-distribution probe agrees: PSM resolves the
interaction variable far better than FB (AUC $0.94$ vs.\ $0.75$), yet scores lower.
(Part of IWR's drop is an inference/measure mismatch --- $B$ whitened under the
reweighted $\rho'$ but the task inferred under $\rho$; the matched GLS inference rule
recovers $+13$--$20$ points, but still below plain FB.) \\
\textbf{Implies:} inside the FB family there is \B{no value-consumer} for interaction
information. Improving the representation alone only perturbs the working
behavior-cloned position policy.

\subsection{G4: success-level objective isolation --- a dissociation}
\textbf{Did:} the inference result above is value-level; the bridge's $94$ uses the
\emph{exact goal} fed to the network. G4 closes the loop at the \emph{policy} level
under FB's \emph{full zero-shot pipeline}: bilinear IQL critic $+$ a flow actor
conditioned on $z{=}\text{stop\_grad}(\psi(g))$ (single shared $B$; stop-grad keeps
the representation purely IQL-trained), 9 seeds @1M, evaluated two ways --- exact
$z{=}\psi(g)$ and zero-shot $z{=}\mathbb{E}[r\,\psi(s')]$ (FB's L0 rule). \\
\textbf{Found:} a sharp \B{dissociation} ($n{=}9$; success with $95\%$ CI):
\begin{center}
\begin{tabular}{lcc}
\toprule
representation objective & actor conditioning & success ($\pm$\,95\% CI) \\
\midrule
measure (plain FB)      & zero-shot $z$ & $48 \pm 5$ \\
measure (FB-IQL critic) & zero-shot $z$ & $24 \pm 7$ \\
\B{IQL (G4)}            & \B{exact $z=\psi(g)$} & $\mathbf{93 \pm 2}$ \\
\B{IQL (G4)}            & \B{zero-shot $z=\mathbb{E}[r\psi(s')]$} & $\mathbf{7 \pm 3}$ \\
IQL (bridge)            & exact goal (raw) & $94 \pm 1$ \\
\bottomrule
\end{tabular}
\end{center}
\textbf{Implies:} with the \emph{exact} goal embedding, the IQL-trained representation
$+$ $z$-conditioned actor essentially matches the raw-goal bridge ($93$ vs.\ $94$; at
1M the cost of routing the goal through the low-rank $\psi$ bottleneck and conditioning
the actor on $z$ rather than $g$ has all but vanished --- it was $\sim9$ pts at 500k).
So the objective, form, rank, and
$z$-conditioning all \emph{compose} at the policy level. But FB's zero-shot inference
rule \B{collapses the policy} ($6.7$), even though the same rule leaves the
\emph{value's ranking} intact (\cfauc{} $0.92$, \S\,inference above). This is the
predicted ``success-level conditioning effect invisible to value-ranking'': a value
that ranks goal-reaching states well under an inferred $z$ does not yield a policy
that \emph{acts} well under it. The breaking point of FB's pipeline at the success
level is the inference\,$\to$\,actor path, not the objective, form, or rank.

\paragraph{Caveat on the inferred-$z$ collapse.} The actor was trained only on
exact-goal embeddings $z=\psi(g)$, so at eval the inferred
$z=\mathbb{E}[r\psi(s')]$ is \emph{out of the actor's training $z$-distribution}; the
relabel batches were also sparse (positive fraction $\approx2\%$). The $6.7$ therefore
conflates ``FB inference is fragile'' with ``this actor never saw inferred-style
$z$.'' A clean attribution needs an actor trained across the inferred-$z$ distribution
(as canonical FB does, sampling $z$ from random-reward readouts) --- the natural next
control. What G4 \emph{does} establish cleanly is the exact-$z$ result: the value
objective suffices end-to-end at the policy level when the task vector is in
distribution.

\section{Cross-domain controls: does the actor/critic story transfer?}
The morphing grid above is entirely on \texttt{cube-single} (state). The natural
control is whether the same knobs behave the same way on other OGBench domains.
We re-ran the key critic$\times$actor cells on \textbf{AntMaze-medium}
(locomotion-navigation, no object) and \textbf{Scene} (multi-object manipulation).
These are shorter, lighter controls than the main grid --- $n{=}3$ seeds at
\textbf{500k} steps (vs.\ $n{=}8$--$10$ at 1M) --- so the intervals are wide;
$\pm$ is the 95\% CI (Student-$t$, $n{-}1$ df). The headline is a
\textbf{domain-dependent reversal}: the flow actor that is \emph{best} on cube is
\emph{worst} on AntMaze.

\subsection{AntMaze-medium}
The full $2\times2$ of critic form (MLP vs.\ bilinear $d{=}50$) $\times$ actor
(Gaussian DDPG+BC vs.\ flow), plus FB, on \texttt{antmaze-medium-navigate-v0}:
\begin{center}
\small
\begin{tabular}{lcc}
\toprule
\textbf{Critic $\times$ actor} & \textbf{Success @500k (95\% CI)} & \textbf{$n$} \\
\midrule
bilinear $d{=}50$ $+$ \textbf{DDPG+BC} actor & $\mathbf{83.7 \pm 2.1}$ & 3 \\
MLP (stock GCIQL) $+$ DDPG+BC actor          & $69.5 \pm 13.4$ & 3 \\
MLP $+$ \textbf{flow} actor                  & $33.3 \pm 16.0$ & 3 \\
bilinear $d{=}50$ $+$ \textbf{flow} actor    & $23.9 \pm 4.9$  & 3 \\
\midrule
FB (flow actor, measure-matching)            & $59.3 \pm 25.0$ & 3 \\
\bottomrule
\end{tabular}
\end{center}
\noindent\textbf{Reading it.} On cube the flow actor is what \emph{realizes} the
bilinear value (bilinear$+$flow $92.6$, bilinear$+$Gaussian $56.6$). On AntMaze the
ordering \emph{flips}: the flow actor is the worst choice under \emph{either} critic
(bilinear $24$, MLP $33$), while DDPG+BC is best under both (bilinear $84$, MLP $70$).
So the cube finding ``the flow actor is needed and held fixed, hence not the cause''
does not generalize --- on navigation the actor parametrization \emph{is} a
first-order lever, in the opposite direction. FB ($59$) sits between the two flow
cells and the DDPG+BC cells, consistent with its flow actor being a real handicap
here rather than the measure objective alone.

\subsection{Scene}
Multi-object manipulation (\texttt{scene-play-v0}), the cells we ran:
\begin{center}
\small
\begin{tabular}{lcc}
\toprule
\textbf{Critic $\times$ actor} & \textbf{Success @500k (95\% CI)} & \textbf{$n$} \\
\midrule
MLP (stock GCIQL) $+$ DDPG+BC actor       & $\mathbf{51.6 \pm 8.9}$ & 3 \\
bilinear $d{=}50$ $+$ flow actor          & $29.2 \pm 7.9$ & 3 \\
\midrule
FB (flow actor, measure-matching)         & $4.7 \pm 5.7$  & 3 \\
\bottomrule
\end{tabular}
\end{center}
\noindent\textbf{Reading it.} Scene echoes AntMaze on the actor axis (bilinear$+$flow
$29 <$ stock DDPG+BC $52$) and echoes \emph{cube} on the FB axis: FB collapses to
near-zero ($4.7$), its worst showing on any domain here. The successor-measure
objective with a flow actor degrades hardest exactly where the interaction variable
is richest (more objects), consistent with the under-resolution story that motivates
the parent draft, while the IQL-trained GCIQL critic stays well above water.

\section{What this says about what could be going on}

Stacking the eliminations:
\begin{enumerate}[nosep]
  \item \B{Actor} --- exonerated (FB's flow actor on a good critic: $73\pm4$).
  \item \B{Bilinear form} --- exonerated (bilinear GCIQL: 94).
  \item \B{Rank} --- exonerated ($d{=}50$: 92.6).
  \item \B{Zero-shot inference} --- exonerated at the value level (\cfauc{}
    $0.92\approx0.92$).
  \item \B{Value-target form} --- exonerated (a correct IQL value on FB's bilinear:
    still 24).
\end{enumerate}

\noindent By elimination, the binding ingredient is the \B{learning objective}:
end-to-end IQL value-TD versus measure-matching. The two directions agree:
\begin{itemize}[nosep]
  \item \emph{IQL\,$\to$\,FB}: pour IQL value-TD through FB's exact form and rank
    (exact goal) and you get $\sim94$ ($93$ once the actor is $z$-conditioned).
  \item \emph{FB\,$\to$\,IQL}: pour the IQL value-target onto FB's bilinear and you
    get $24$ \emph{at the old incidentals} (lr 1e-4, ortho 1000, cube-xyz reward) ---
    \textcolor{red}{but $\sim$59 at sane incidentals (lr 3e-4, ortho 0, goal-matching
    reward; Part~1 ladder)}, so this direction is \emph{not} a flat value-blind wall:
    the objective lifts FB well above the future-reachability baselines, it just stops
    short of GCIQL ($59\to90$). Making it a \emph{better measure} representation (IWR,
    PSM, $d{=}128$) still makes success \emph{worse}.
\end{itemize}

\noindent \B{The proposed mechanism.} Measure-matching objectives shape the
representation to reconstruct the discounted state-occupancy density. On manipulation
this density is dominated by where the \emph{arm} goes and by the gross geometry of
the cube; the thin, interaction-conditioned successor structure that decides whether
a grasp will be \emph{kept} through transport is a low-mass correction the density
objective has no incentive to resolve into the value. A value-TD objective, by
contrast, propagates the goal-reaching signal backward through exactly those
transitions, so the value becomes sensitive to the keep-vs-drop distinction even in
the same bilinear, low-rank form. The fix is therefore not a better basis but an
\B{end-to-end value-TD coupling} of representation, value, and task.

\noindent \B{What G4 adds.} The success-level test confirms this end-to-end at the
policy level when the task vector is \emph{in distribution}: an IQL-trained bilinear
$+$ $z$-conditioned actor reaches $93$ with the exact goal embedding (vs.\ FB's $48$
and FB-IQL's $24$). But it also exposes a \emph{second, separable} failure mode that
the value-level probes missed: under FB's \emph{zero-shot} inference the policy
collapses to $7$ even though the value still ranks correctly (\cfauc{} $0.92$). So the
full story has two layers --- (i) the \B{objective} sets whether the representation's
value is controllable at all (the $48\!\to\!93$ jump), and (ii) the
\B{inference\,$\to$\,actor} path is an independent success-level bottleneck once you
go zero-shot. Value-ranking robustness to inferred $z$ does not transfer to policy
robustness. The clean next control is an actor trained across the inferred-$z$
distribution, to separate ``FB inference is fragile'' from ``this actor was never
trained on inferred $z$.''

\paragraph{Caveats.} The starting-point table (\S1) is each method's \emph{best} eval over
steps $\ge$500k (peak step varies by method; the $\ge$500k floor drops early transients
like CRL's 200k spike). \emph{The ablation grid and FB-side interventions are now
reported at 1M} ($n{=}8$--$10$ seeds; the 10-seed\,$\times$\,1M refresh, campaign spec
\texttt{2026-06-22-bridge-table-10seed-1M}); a few auxiliary rows (IWR, FB $d{=}128$,
PSM) remain at their original snapshot and are flagged inline. The 1M grid is what the
bridge controls against --- e.g.\ GCIQL is $78$ at its peak but $73$ at 1M. $n{=}10$
for most rows ($n{=}8$ for $d{=}512$, $n{=}9$ for FlowBC/G4/CRL-flow), $n{=}3$ for the
few remaining auxiliary rows (mean, not IQM, with a correspondingly wide $\pm$).
\cfauc{}/lift are value-level probes on the bilinear
checkpoints (representation properties, step-robust); G4 closes the success-level /
on-$\rho$ arm. Faithful reward uses privileged cube qpos (the observation cube xy has
unrecoverable per-axis scale/offset; the offline dataset lacks qpos).

\end{document}
